\section{50-State Reference Table}
\label{sec:fifty_state_table}

Table~\ref{tab:fifty_state} provides the definitive 50-state reference for population definition sensitivity. For each state, we report: the number of congressional seats, the total-vs.-CVAP tract reassignment rate, the total-vs.-prison-adjusted tract reassignment rate, the partisan direction (D = Democratic-advantaging switch, R = Republican-advantaging switch, N = neutral $<$0.3pp), and the RSI change under CVAP redistricting.

States are grouped by degree of CVAP effect: high ($\geq$5\% reassignment), moderate (1--5\%), and low ($<$1\%).

\begin{longtable}{llccccc}
\caption{50-State Population Definition Sensitivity Reference Table. ``CVAP RR'' = tract reassignment rate switching from total to citizen VAP. ``Adj RR'' = tract reassignment rate switching from total to prison-adjusted total. ``Dir-CVAP'' = partisan direction of CVAP switch (R=Republican, D=Democratic, N=neutral). ``RSI$\Delta$'' = RSI change under CVAP redistricting.} \label{tab:fifty_state} \\
\toprule
\textbf{State} & \textbf{Seats} & \textbf{CVAP RR} & \textbf{Adj RR} & \textbf{Dir-CVAP} & \textbf{Dir-Adj} & \textbf{RSI$\Delta$} \\
\midrule
\endfirsthead
\multicolumn{7}{c}{\tablename\ \thetable{} -- continued} \\
\toprule
\textbf{State} & \textbf{Seats} & \textbf{CVAP RR} & \textbf{Adj RR} & \textbf{Dir-CVAP} & \textbf{Dir-Adj} & \textbf{RSI$\Delta$} \\
\midrule
\endhead
\midrule
\multicolumn{7}{r}{\textit{Continued on next page}} \\
\endfoot
\bottomrule
\endlastfoot
%
% HIGH EFFECT STATES (CVAP >= 5%)
\multicolumn{7}{l}{\textit{High effect: CVAP reassignment $\geq$ 5\%}} \\
California & 52 & 11.8\% & 1.0\% & R & D & $-1.4\%$ \\
Texas & 38 & 9.4\% & 2.1\% & R & D & $-0.9\%$ \\
New York & 26 & 8.2\% & 1.2\% & R & D & $-0.8\%$ \\
Florida & 28 & 6.7\% & 0.6\% & R & D & $-0.7\%$ \\
Nevada & 4 & 5.8\% & 0.3\% & R & N & $-2.1\%$ \\
Arizona & 9 & 5.3\% & 0.4\% & R & N & $-0.5\%$ \\
New Jersey & 12 & 5.1\% & 0.5\% & R & D & $-0.6\%$ \\
Hawaii & 2 & 5.0\% & 0.1\% & R & N & $-2.3\%$ \\
%
% MODERATE EFFECT STATES (1% <= CVAP < 5%)
\multicolumn{7}{l}{\textit{Moderate effect: CVAP reassignment 1--5\%}} \\
Illinois & 17 & 4.9\% & 1.6\% & R & D & $-0.5\%$ \\
Washington & 10 & 3.4\% & 0.4\% & R & N & $-0.3\%$ \\
Virginia & 11 & 2.8\% & 0.6\% & R & D & $-0.4\%$ \\
Georgia & 14 & 2.6\% & 0.7\% & R & D & $-0.3\%$ \\
Colorado & 8 & 2.5\% & 0.3\% & R & N & $-0.3\%$ \\
Maryland & 8 & 2.4\% & 0.4\% & R & D & $-0.4\%$ \\
North Carolina & 14 & 2.2\% & 0.8\% & R & D & $-0.3\%$ \\
Massachusetts & 9 & 2.0\% & 0.3\% & N & N & $-0.2\%$ \\
Connecticut & 5 & 1.9\% & 0.4\% & N & N & $-0.2\%$ \\
Minnesota & 8 & 1.8\% & 0.6\% & N & D & $-0.2\%$ \\
Oregon & 6 & 1.7\% & 0.3\% & R & N & $-0.2\%$ \\
Utah & 4 & 1.6\% & 0.2\% & R & N & $-0.2\%$ \\
Kansas & 4 & 1.5\% & 0.5\% & R & N & $-0.1\%$ \\
Michigan & 13 & 1.5\% & 0.7\% & R & D & $-0.2\%$ \\
Pennsylvania & 17 & 1.4\% & 1.8\% & N & D & $-0.2\%$ \\
Nebraska & 3 & 1.3\% & 0.3\% & R & N & $-0.1\%$ \\
Rhode Island & 2 & 1.3\% & 0.2\% & N & N & $-0.2\%$ \\
Wisconsin & 8 & 1.2\% & 0.6\% & N & D & $-0.1\%$ \\
Delaware & 1 & 1.2\% & 0.3\% & N & N & --- \\
Oklahoma & 5 & 1.1\% & 0.7\% & N & D & $-0.1\%$ \\
Ohio & 15 & 1.1\% & 1.4\% & N & D & $-0.1\%$ \\
Indiana & 9 & 1.0\% & 0.8\% & N & D & $-0.1\%$ \\
%
% LOW EFFECT STATES (CVAP < 1%)
\multicolumn{7}{l}{\textit{Low effect: CVAP reassignment $<$ 1\%}} \\
Tennessee & 9 & 0.9\% & 0.9\% & N & D & $0.0\%$ \\
Alabama & 7 & 0.8\% & 1.1\% & N & D & $0.0\%$ \\
Missouri & 8 & 0.8\% & 0.7\% & N & D & $0.0\%$ \\
Arkansas & 4 & 0.7\% & 0.8\% & N & D & $0.0\%$ \\
Iowa & 4 & 0.7\% & 0.5\% & N & N & $0.0\%$ \\
Louisiana & 6 & 0.7\% & 0.9\% & N & D & $0.0\%$ \\
South Carolina & 7 & 0.6\% & 0.7\% & N & D & $0.0\%$ \\
New Mexico & 3 & 0.6\% & 0.3\% & N & N & $0.0\%$ \\
Idaho & 2 & 0.5\% & 0.2\% & N & N & $0.0\%$ \\
Mississippi & 4 & 0.5\% & 1.0\% & N & D & $0.0\%$ \\
Kentucky & 6 & 0.5\% & 0.9\% & N & D & $0.0\%$ \\
New Hampshire & 2 & 0.4\% & 0.2\% & N & N & $0.0\%$ \\
Maine & 2 & 0.4\% & 0.2\% & N & N & $0.0\%$ \\
West Virginia & 2 & 0.3\% & 0.5\% & N & N & $0.0\%$ \\
Vermont & 1 & 0.3\% & 0.1\% & N & N & --- \\
Alaska & 1 & 0.3\% & 0.2\% & N & N & --- \\
North Dakota & 1 & 0.2\% & 0.1\% & N & N & --- \\
South Dakota & 1 & 0.2\% & 0.1\% & N & N & --- \\
Montana & 2 & 0.2\% & 0.1\% & N & N & $0.0\%$ \\
Wyoming & 1 & 0.1\% & 0.1\% & N & N & --- \\
\midrule
\textbf{National (wtd)} & \textbf{435} & \textbf{3.2\%} & \textbf{0.8\%} & \textbf{R} & \textbf{D} & \textbf{$-0.5\%$} \\
\end{longtable}

\noindent\textbf{Notes:} CVAP RR and Adj RR are population-weighted tract reassignment rates within each state. ``Dir-CVAP'' and ``Dir-Adj'' show the partisan direction of the switch among affected districts: R = Republican-advantaging ($\Delta L < -0.3$pp), D = Democratic-advantaging ($\Delta L > +0.3$pp), N = neutral ($|\Delta L| \leq 0.3$pp). RSI$\Delta$ is the mean change in Reock Shape Index across all state congressional districts; a negative value indicates a slight decrease in compactness (less compact on average). ``---'' indicates single-seat states where within-state redistricting does not apply. CVAP data from 2019 ACS 5-year estimates; adjustment data from 2020 Census Group Quarters and PPI 2020 facility database.
